\chapter{Modulo de Código.}

Para inicializar un nuevo proyecto con ejemplos de java. Esto creará un proyecto en la carpeta examples\_project

\begin{lstlisting}[style=terminal]
# En el nuevo proyecto de notas, desde la raíz:
mkdir -p rsc/code/examples_project
cd rsc/code/examples_project
gradle init --type java-application --dsl groovy --package com.examples
# Eliminar el .git en examples_project por que puede causar problemas
rm -rf rsc/code/examples_project/.git
# Luego editar build.gradle y settings.gradle con tus parámetros
\end{lstlisting}

Después de crear el proyecto se puede abrir con Intellij la carpeta examples\_project.

Los parámetros a editar en el archivo \code{build.gradle} son: 

\begin{itemize}
    \item Group
    \item Version
    \item Java version
    \item Dependencies
    \item mainClass
\end{itemize}

\lstinputlisting[
    caption={Archivo de gradle con parámetros configurados.}, 
    label={lst:build.gradle file},
    firstline=9, lastline=43, 
    language=Java]
{rsc/code/examples_project/app/build.gradle}

En este punto el proyecto compila correctamente y se pueden agregar ejemplos creando paquetes bajo com.examples (e.g. com.examples.ch01.streams, com.examples.ch02.collections).

\section{Resaltado de código.}

    \subsection{Código en línea.}
    
    Se utiliza para indicar fragmentos de código dentro del texto.


    \nkCode{source ~/.ansible-venv/bin/activate}
    
    \nkCode{git checkout -- rsc/pckg/template-package/noteskit-sections.sty}
    
    \nkCode{string name = "Mau";}
    
    \nkCode{ssh-keyscan 192.168.56.11 >> ~/.ssh/known_hosts}
    
    \code{ssh-keyscan 192.168.56.11 >> ~/.ssh/known_hosts}

    \subsection{Comandos en línea.}
    
    Se utiliza para indicar comandos de la CLI dentro del texto o una tabla.


    \nkCmd{HashMap<String, ~Int>}
    
    
    \nkCmd{source ~/.ansible-venv/bin/activate}
    
    \nkCmd{git checkout -- rsc/pckg/template-package/noteskit-sections.sty}
    
    \nkCmd{string name = "Mau";}
    
    \nkCmd{ssh-keyscan 192.168.56.11 >> ~/.ssh/known_hosts}
    
    \cmd{ssh-keyscan 192.168.56.11 >> ~/.ssh/known_hosts}
    
    \begin{nkTable}{|X[l,m,1]|X[c,m,10]|X[l,m,40]|}{Ejemplo de tabla con comandos.}{tab:command_table_example}
        \textbf{\#} &	Instrucción & Descripción  \\
        1 & \nkCmd{python -m vent ~/.ansible-venv} & Crea un entorno virtual de Python en \nkCode{~/.ansible-venv}.\\
    \end{nkTable}
    
    \begin{nkTable}{|c|c|X|}{Ejemplo de tabla.}{tab:command_table_example2}
        \textbf{\#} & \textbf{Instrucción} & \textbf{Descripción} \\ 
        1 & \nkCmd{python -m vent ~/.ansible-venv} &  Crea un entorno virtual de Python en \nkCode{/.ansible-venv}.\\
        2 &  & Software para ejecutar programas Java. \\    
    \end{nkTable}
    
    %\lstinline[style=terminal]{~/.ansible-venv}
    
    \begin{tblr}{colspec={|c|c|c|}}
        {\ttfamily\small\color{OliveGreen}\fakeverb{~/.ansible-venv}}& $a\,b$ & \\
    \end{tblr}
    
    otra tabla 
    
    \begin{tblr}{colspec={|X|X|X|}}
        {\ttfamily\small\color{OliveGreen}\fakeverb{~/.ansible-venv}}& $a\,b$ & \\
    \end{tblr}
    
    \subsection{Bloques de código.}
    
    \begin{lstlisting}[style=IntelliJLight]
/**
* @author John Doe
* @param args Command line arguments
*/
@Override
public static void main(String[] args) {
    String message = "Hello World";
    // TODO: Add more functionality
    System.out.println(message);
}
    \end{lstlisting}
    
    \lstinputlisting[style=IntelliJDark, caption={A listing from the file
        \ttfamily{Rectangle.java}}, label={lst:file}]{rsc/code/Rectangle.java}

